% QUESTION: 7
% STATUS: DRAFT
% AUTHOR: DG
% CREATED: 2026-02-02
% LAST UPDATED: 2026-02-09
% NOTES: 
% Q: If we perform the exact same operation twice, do we need to explain it in the same level of detail the second time? For example, if we showed how to apply L’Hôpital’s Rule in detail in Step 2, do we need to repeat the full explanation again in Step 4 when we apply L’Hôpital’s Rule a second time?
% Step 4 needs to add more steps.
% Maybe split Step 2 into 2 big steps. Right now it has 12 substeps, which is too long, and the substeps need to be further expand.

% Question
\section*{Question 7}
Evaluate $\displaystyle\lim_{x\to 0}\frac{1-\cos 2x}{x^3+5x^2}$.
    
% Multiple Choices
\begin{enumerate}[label = \alph*.]
    \item 0
    \item $\displaystyle\frac{1}{10}$
    \item $\displaystyle\frac{2}{5}$
    \item I don't know
\end{enumerate}

\setcounter{solutionMainStep}{0}
\setcounter{solutionSubStep}{0}

% Solution Box
\begin{tcolorbox}[breakable, colback= blue!5!white, colframe= blue!50!black, title=\textbf{Solution: c) $\displaystyle\frac{2}{5}$}]

    % Step 1
    \mainStep{7}{Identify the limit form}

    % Step 1.1
    \solutionStep{7}{Evaluate numerator at $x=0$}{
    \begin{flushleft}
    \begin{tabbing}
    \hspace{4.5cm} \= \kill
    $\begin{aligned}[t]
    1-\cos 2x &= 1-\cos(2\cdot0) \\
              &= 1-\cos(0) \\
              &= 1-1 \\
              &= 0
    \end{aligned}$
    \end{tabbing}
    \end{flushleft}
    }

    % Step 1.2
    \solutionStep{7}{Evaluate denominator at $x=0$}{
    \begin{flushleft}
    \begin{tabbing}
    \hspace{4.5cm} \= \kill
    $\begin{aligned}[t]
    x^3+5x^2 &= (0)^3+5(0)^2 \\
             &= 0+0 \\
             &= 0
    \end{aligned}$
    \end{tabbing}
    \end{flushleft}
    }

    % Step 1.3
    \solutionStep{7}{Determine indeterminate form}{
    $\lim_{x\to 0}\frac{1-\cos 2x}{x^3+5x^2} = \frac{0}{0}$}

    % Step 2
    \mainStep{7}{Apply L'Hôpital's Rule (First Time)}

    % Step 2.1
    \solutionStep{7}{State L'Hôpital's Rule}{
    \colorbox{yellow}{L'Hôpital's Rule:} If $\lim_{x\to a}\frac{f(x)}{g(x)}$ gives $\frac{0}{0}$ or $\frac{\infty}{\infty}$, then 
    $\displaystyle\lim_{x\to a}\frac{f(x)}{g(x)}=\lim_{x\to a}\frac{f'(x)}{g'(x)}$, provided the new limit exists.
    }

    % Step 2.2
    \solutionStep{7}{Differentiate numerator}{
    $\frac{d}{dx}(1-\cos 2x)
    =
    \frac{d}{dx}(1)
    +
    \frac{d}{dx}(-\cos 2x)$
    }

    % Step 2.3
    \solutionStep{7}{Differentiate $1$}{
    $\frac{d}{dx}(1)=0$
    }

    % Step 2.4
    \solutionStep{7}{Differentiate $-\cos 2x$ using chain rule}{
    \colorbox{yellow}{Chain Rule:} If $F(x)=\cos(h(x))$, then $\frac{d}{dx}[\cos(h(x))]=-\sin(h(x))h'(x)$.
    }

    % Step 2.5
    \solutionStep{7}{Differentiate inside function}{
    $\frac{d}{dx}(2x)=2$
    }

    % Step 2.6
    \solutionStep{7}{Combine chain rule result}{
    $\frac{d}{dx}(-\cos 2x)=(-1)(-\sin 2x)(2)=2\sin 2x$
    }

    % Step 2.7
    \solutionStep{7}{Combine numerator derivatives}{
    $\frac{d}{dx}(1-\cos 2x)=0+2\sin 2x=2\sin 2x$
    }

    % Step 2.8
    \solutionStep{7}{Differentiate denominator}{
    $\frac{d}{dx}(x^3+5x^2)
    =
    \frac{d}{dx}(x^3)
    +
    \frac{d}{dx}(5x^2)$
    }

    % Step 2.9
    \solutionStep{7}{Apply power rule}{
    \colorbox{yellow}{Power Rule:} If $g(x)=x^n$, then $g'(x)=nx^{n-1}$.
    }

    % Step 2.10
    \solutionStep{7}{Differentiate $x^3$}{
    $\frac{d}{dx}(x^3)=3x^2$
    }

    % Step 2.11
    \solutionStep{7}{Differentiate $5x^2$}{
    $\frac{d}{dx}(5x^2)=10x$
    }

    % Step 2.12
    \solutionStep{7}{Combine denominator derivatives}{
    $\frac{d}{dx}(x^3+5x^2)=3x^2+10x$
    }

    % Step 2.13
    \solutionStep{7}{Rewrite limit}{
    $\displaystyle\lim_{x\to0}\frac{2\sin 2x}{3x^2+10x}$
    }

    % Step 3
    \mainStep{7}{Check the new limit form}

    % Step 3.1
    \solutionStep{7}{Evaluate numerator at $x=0$}{
    \begin{flushleft}
    \begin{tabbing}
    \hspace{4.5cm} \= \kill
    $\begin{aligned}[t]
    2\sin 2x &= 2\sin(2\cdot0) \\
             &= 2\sin(0) \\
             &= 2(0) \\
             &= 0
    \end{aligned}$
    \end{tabbing}
    \end{flushleft}
    }

    % Step 3.2
    \solutionStep{7}{Evaluate denominator at $x=0$}{
    \begin{flushleft}
    \begin{tabbing}
    \hspace{4.5cm} \= \kill
    $\begin{aligned}[t]
    3x^2+10x &= 3(0)^2+10(0) \\
             &= 0+0 \\
             &= 0
    \end{aligned}$
    \end{tabbing}
    \end{flushleft}
    }

    % Step 3.3
    \solutionStep{7}{Determine indeterminate form}{
    $\lim_{x\to0}\frac{2\sin 2x}{3x^2+10x} = \frac{0}{0}$
    }

    % Step 4
    \mainStep{7}{Apply L'Hôpital's Rule (Second Time)}

    % Step 4.1
    \solutionStep{7}{Differentiate numerator}{
    $\frac{d}{dx}(2\sin 2x)=2(\cos 2x)(2)=4\cos 2x$
    }

    % Step 4.2
    \solutionStep{7}{Differentiate denominator}{
    $\frac{d}{dx}(3x^2+10x)=6x+10$
    }

    % Step 4.3
    \solutionStep{7}{Rewrite limit}{
    $\displaystyle\lim_{x\to0}\frac{4\cos 2x}{6x+10}$
    }

    % Step 5
    \mainStep{7}{Evaluate the limit}

    % Step 5.1
    \solutionStep{7}{Evaluate numerator at $x=0$}{
    $4\cos 2x=4\cos(2\cdot0)=4\cos(0)=4(1)=4$
    }

    % Step 5.2
    \solutionStep{7}{Evaluate denominator at $x=0$}{
    $6x+10=6(0)+10=10$
    }

    % Step 5.3
    \solutionStep{7}{Combine the result}{
    $\lim_{x\to0}\frac{4\cos 2x}{6x+10}=\frac{4}{10}=\frac{2}{5}$
    }

    % Step 6
    \mainStep{7}{Final Answer: the correct option is C. $\displaystyle\frac{2}{5}$ }

\end{tcolorbox}
\newpage
